\documentclass[aspectratio=169,10pt]{beamer}

% Shared Beamer setup for Fall 2026 course slides.
% Clean Cal Poly-inspired theme: no custom class files or uncommon packages.
\mode<presentation>{
  \usetheme{default}
  \usefonttheme{professionalfonts}
}

\definecolor{CalPolyGreen}{HTML}{154734}
\definecolor{CalPolyGold}{HTML}{BD8B13}
\definecolor{CalPolySage}{HTML}{7FA28F}
\definecolor{CalPolyMint}{HTML}{DDF1E7}
\definecolor{CalPolySand}{HTML}{A29F86}
\definecolor{CalPolyBeige}{HTML}{EFEEDE}
\definecolor{DeckInk}{HTML}{1F2933}
\definecolor{DeckMuted}{HTML}{5F6B7A}
\definecolor{DeckSoft}{HTML}{F7F7F5}

\setbeamersize{text margin left=0.55in,text margin right=0.55in}
\setbeamercolor{normal text}{fg=DeckInk,bg=white}
\setbeamercolor{structure}{fg=CalPolyGreen}
\setbeamercolor{title}{fg=white}
\setbeamercolor{frametitle}{fg=CalPolyGreen,bg=white}
\setbeamercolor{block title}{bg=CalPolySage,fg=white}
\setbeamercolor{block body}{bg=CalPolyMint,fg=black}
\setbeamercolor{alerted text}{fg=CalPolyGold}

\setbeamerfont{title}{size=\Huge,series=\bfseries}
\setbeamerfont{frametitle}{size=\Large,series=\bfseries}
\setbeamerfont{block title}{size=\normalsize,series=\bfseries}

\setbeamertemplate{navigation symbols}{}
\setbeamertemplate{itemize item}{\color{CalPolyGold}$\blacktriangleright$}
\setbeamertemplate{itemize subitem}{\color{CalPolyGreen}$\bullet$}
\setbeamertemplate{footline}{
  \hfill{\color{DeckMuted}\scriptsize\insertframenumber/\inserttotalframenumber}
  \hspace{0.18in}\vspace{0.12in}
}
\setbeamertemplate{frametitle}{
  \vspace{0.18in}
  \hspace{0.02in}\insertframetitle\par
  \vspace{0.08in}
  \noindent\makebox[\linewidth][l]{%
    {\color{CalPolyGold}\rule{0.55in}{2pt}}%
    {\color{CalPolyGreen}\rule{\dimexpr\linewidth-0.55in\relax}{2pt}}%
  }
}

\newcommand{\CourseNumber}{Course}
\newcommand{\CourseTitle}{Course Title}
\newcommand{\LectureNumber}{0}
\newcommand{\LectureTitle}{Lecture Title}
\newcommand{\LectureDate}{Date}
\newcommand{\CourseUnit}{Unit}

\newcommand{\coursenumber}[1]{\renewcommand{\CourseNumber}{#1}}
\newcommand{\coursetitle}[1]{\renewcommand{\CourseTitle}{#1}}
\newcommand{\lecturenumber}[1]{\renewcommand{\LectureNumber}{#1}}
\newcommand{\lecturetitle}[1]{\renewcommand{\LectureTitle}{#1}}
\newcommand{\lecturedate}[1]{\renewcommand{\LectureDate}{#1}}
\newcommand{\courseunit}[1]{\renewcommand{\CourseUnit}{#1}}

\author{Kirk Duran}
\institute{California Polytechnic State University, San Luis Obispo}
\date{\LectureDate}

\newcommand{\makedecktitle}{
  \begingroup
  \setbeamercolor{background canvas}{bg=CalPolyGreen}
  \begin{frame}[plain]
    \vfill
    \begin{center}
      {\usebeamerfont{title}\color{white}\LectureTitle\par}
      \vspace{0.18in}
      {\Large\color{white}Lecture \LectureNumber\par}
    \end{center}
    \vfill
    \noindent{\color{white}\rule{\linewidth}{1.2pt}}\par
    \vspace{0.08in}
    \noindent\colorbox{CalPolyGold}{\makebox[0.24\linewidth][c]{\color{white}\Large\CourseNumber}}
    \hspace{0.12in}{\color{white}\Large Kirk Duran}
  \end{frame}
  \endgroup
}

\newenvironment{keyidea}{\begin{block}{Key idea}}{\end{block}}
\newenvironment{activity}{\begin{block}{In class}}{\end{block}}
\newenvironment{cpdefinition}{\begin{block}{Definition}}{\end{block}}
\newenvironment{exampleblockcp}{
  \begingroup
  \setbeamercolor{block title}{bg=CalPolySand,fg=white}
  \setbeamercolor{block body}{bg=CalPolyBeige,fg=black}
  \begin{block}{Example}
}{
  \end{block}
  \endgroup
}

\newcommand{\imageplaceholder}[2][0.3in]{%
  \noindent\makebox[\linewidth][c]{%
    \fcolorbox{CalPolySage}{CalPolyBeige}{%
      \parbox[c][#1][c]{0.86\linewidth}{%
        \centering
        {\color{DeckMuted}\scriptsize Image placeholder: }%
        {\color{DeckInk}\scriptsize\ttfamily\detokenize{#2}}%
      }%
    }%
  }%
}

\newcommand{\sectionframe}[1]{
  \begingroup
  \setbeamercolor{background canvas}{bg=CalPolyGreen}
  \begin{frame}[plain]
    \vfill
    {\color{white}\LARGE\bfseries #1\par}
    \vspace{0.18in}
    {\color{CalPolyGold}\rule{0.22\paperwidth}{3pt}}
    {\color{white}\rule{0.62\paperwidth}{3pt}}
    \vfill
  \end{frame}
  \endgroup
}

\newcommand{\placeholdernote}{
  \begin{block}{Placeholder}
    Replace these bullets with the final lecture material, examples, and in-class activities.
  \end{block}
}

\usepackage[T1]{fontenc}
\usepackage{lmodern}
\usepackage{amsmath,amssymb}
\usepackage{graphicx}
\usepackage{booktabs}
\usepackage{array}
\pdfstringdefDisableCommands{\def\translate#1{#1}}
\graphicspath{{../assets/lecture-32/}}

% Cal Poly CS 4880, Lecture 32, November 13, 2026.
% Rebuilt after reading all 17 slides of Lecture 22, Learning - Reinforcement Learning.
% About 50 minutes. Original vacuum example, +100 completion/-10 action;
% separate simplified 5x5 GridWorld example, +100 goal/-1 move.
% Illustrations and reported measurements are imported from the supplied image archive.
% The original figure generator was not supplied; reported run uses seed 4880.

\newcommand{\lecturevisual}[2][2.15in]{%
  \par\vspace{2pt}{\centering\includegraphics[width=.98\linewidth,height=#1,keepaspectratio]{#2}\par}\vspace{2pt}%
}
\newcommand{\architecturestage}[1]{%
  {\centering\setlength{\fboxsep}{5pt}%
  \fcolorbox{CalPolyGreen}{DeckSoft}{\parbox{.86\linewidth}{\centering\small #1}}\par}%
}
\newcommand{\conceptlead}[1]{{\small\bfseries\color{CalPolyGreen}#1\par}}
\newenvironment{calculation}[1][Mathematical description]{%
  \par\begingroup\setbeamercolor{calculation body}{bg=DeckSoft,fg=DeckInk}%
  \begin{beamercolorbox}[sep=4pt,wd=\linewidth]{calculation body}%
  \footnotesize{\bfseries\color{DeckMuted}#1\par}%
  \setlength{\jot}{0pt}\setlength{\abovedisplayskip}{3pt}\setlength{\belowdisplayskip}{3pt}%
}{\end{beamercolorbox}\endgroup}

\setbeamersize{text margin left=.42in,text margin right=.42in}
\coursenumber{CS 4880}
\coursetitle{Artificial Intelligence}
\lecturenumber{32}
\lecturetitle{Reinforcement Learning: Rewards, Policies, and Q-Learning}
\newcommand{\LectureDisplayTitle}{Reinforcement Learning:\\[3pt]Rewards, Policies,\\[3pt]and Q-Learning}
\lecturedate{November 13, 2026}
\courseunit{Learning Agents and Machine Learning}

\begin{document}
\begingroup
\renewcommand{\LectureTitle}{\LectureDisplayTitle}
\makedecktitle
\endgroup

\section{Learning Through Interaction}

\begin{frame}{Learning Agents Improve Decisions from Outcomes}
\small
\conceptlead{Unlike labeled examples, reinforcement learning supplies feedback after an action changes the world.}
\begin{columns}[T,onlytextwidth]
\begin{column}{.54\textwidth}
\lecturevisual[2.32in]{agent_loop.png}
\end{column}
\begin{column}{.43\textwidth}
\begin{itemize}
\item An intelligent agent observes its environment, chooses actions, and receives feedback.
\item The core cycle is \textbf{observe $\to$ decide $\to$ act $\to$ learn}.
\item Neural networks can approximate perception, decisions, or action values when explicit tables become impractical.
\end{itemize}
\end{column}
\end{columns}
\begin{keyidea}We want a policy that performs well over repeated interactions, not just one correct prediction.\end{keyidea}
\end{frame}

\begin{frame}{Our Running Example: A Vacuum Robot}
\small
\conceptlead{Imagine a robot cleaning a room with obstacles and dirty tiles.}
\begin{columns}[T,onlytextwidth]
\begin{column}{.54\textwidth}
\lecturevisual[2.13in]{vacuum_grid.png}
\end{column}
\begin{column}{.43\textwidth}
\begin{itemize}
\item Dark cells are obstacles; brown cells are dirty; the teal circle is the robot.
\item The robot knows its location, and movement is deterministic in this toy problem.
\item It can move up, down, left, right, or clean.
\item Search could plan every time; learning instead reuses experience.
\end{itemize}
\end{column}
\end{columns}
\end{frame}

\begin{frame}{What Exactly Is a State?}
\footnotesize
\conceptlead{A state must distinguish situations that require different decisions.}
\begin{columns}[T,onlytextwidth]
\begin{column}{.54\textwidth}
\begin{itemize}
\item The vacuum state has the robot's \textbf{position} and the \textbf{clean/dirty status} of tiles.
\item Two visits to the same square may be different states if the remaining dirt differs.
\item A useful state representation captures information needed to predict the next state and reward.
\end{itemize}
\begin{calculation}[Illustrative vacuum state]
\[s=(\text{row},\text{column},\text{dirt configuration}).\]
\end{calculation}
\end{column}
\begin{column}{.43\textwidth}
\lecturevisual[1.73in]{vacuum_grid.png}
\end{column}
\end{columns}
\begin{keyidea}The robot's position alone is not a complete state for the full cleaning task.\end{keyidea}
\end{frame}

\begin{frame}{Actions Cause Transitions}
\small
\conceptlead{The environment maps a current state and action into a next state and reward.}
\lecturevisual[1.52in]{state_transition.png}
\begin{itemize}
\item For the vacuum, actions are movement or cleaning; a move may be blocked by a wall.
\item A \textbf{deterministic} transition always produces the same outcome from the same state-action pair.
\item Real robots often have \textbf{stochastic} motion or incomplete information; we start with the easier case.
\end{itemize}
\end{frame}

\begin{frame}{Reward Turns a Task Into an Optimization Problem}
\small
\conceptlead{Not every successful behavior is equally efficient.}
\begin{columns}[T,onlytextwidth]
\begin{column}{.52\textwidth}
\begin{itemize}
\item \textbf{Original vacuum example:} $+100$ for cleaning the entire room and $-10$ for each action.
\item Completing the task matters, but taking a long detour is costly.
\item These are \emph{designer-chosen} numbers, not quantities learned automatically.
\item If reward is poorly designed, an agent may optimize the wrong behavior.
\end{itemize}
\end{column}
\begin{column}{.45\textwidth}
\lecturevisual[2.31in]{reward_tradeoff.png}
\end{column}
\end{columns}
\begin{keyidea}Reward describes what the agent is encouraged to do; it is not a guarantee of sensible behavior.\end{keyidea}
\end{frame}

\begin{frame}{Immediate Reward Is Not the Whole Objective}
\small
\conceptlead{An action can be temporarily costly but lead to a better future.}
\begin{calculation}[Discounted return after time $t$]
\[G_t=r_{t+1}+\gamma r_{t+2}+\gamma^2r_{t+3}+\cdots
=\sum_{k=0}^{\infty}\gamma^k r_{t+k+1}.\]
The discount factor $0\leq\gamma<1$ reduces the weight of distant rewards in a continuing task.
\end{calculation}
\begin{itemize}
\item A direct route may cost fewer actions than a wandering route, even if both eventually finish.
\item A policy seeks high \textbf{expected return}, not necessarily the largest immediate reward.
\item For finite episodes, the sum stops at termination; $\gamma=1$ can also be appropriate.
\end{itemize}
\end{frame}

\begin{frame}{A Policy Is the Agent's Strategy}
\small
\conceptlead{A policy says what the robot will do when it encounters a state.}
\begin{columns}[T,onlytextwidth]
\begin{column}{.49\textwidth}
\begin{calculation}[Deterministic policy]
\[\pi(s)=a.\]
A mapping from each state $s$ to one selected action $a$.
\end{calculation}
\begin{itemize}
\item A suboptimal policy wastes moves.
\item A random policy ignores what the robot has learned.
\item An optimal policy maximizes expected discounted return under the task model.
\end{itemize}
\end{column}
\begin{column}{.48\textwidth}
\lecturevisual[2.3in]{policy_vs_q.png}
\end{column}
\end{columns}
\begin{keyidea}Rather than plan from scratch every time, the agent can learn which action is best in each state.\end{keyidea}
\end{frame}

\section{Q-Learning: Assign Values to Actions}

\begin{frame}{Q-Learning Asks: How Good Is Each Action Here?}
\small
\conceptlead{Q-learning estimates the value of choosing action $a$ in state $s$ and acting well afterward.}
\begin{calculation}[Optimal action-value function]
\[Q^*(s,a)=\mathbb{E}\!\left[r_{t+1}+\gamma r_{t+2}+\cdots\mid s_t=s,\,a_t=a,\,\text{optimal continuation}\right].\]
\end{calculation}
\begin{itemize}
\item A \textbf{Q-value} summarizes an immediate choice plus its possible consequences.
\item The algorithm can learn from observed transitions without knowing the transition function in advance.
\item After learning, the greedy policy selects the highest-valued action:
\end{itemize}
\begin{calculation}[Greedy action choice]
\[\pi(s)=\underset{a\in A(s)}{\arg\max}\;Q(s,a).\]
\end{calculation}
\end{frame}

\begin{frame}{The Q-Table Is a State-by-Action Lookup Table}
\small
\conceptlead{Rows represent states; columns represent actions; cells hold estimates of long-term value.}
\lecturevisual[1.85in]{qtable_before_after.png}
\begin{itemize}
\item At initialization, entries are often all zero; experience updates the visited pairs.
\item A row's largest entry determines a greedy action; illegal actions are excluded.
\end{itemize}
\end{frame}

\begin{frame}{How One Experience Updates the Table}
\small
\conceptlead{We move the current estimate toward a target built from reward plus the best next-state value.}
\begin{calculation}[Q-learning temporal-difference update]
\[Q(s,a)\leftarrow Q(s,a)+\alpha\underbrace{\left[r+\gamma\max_{a'}Q(s',a')-Q(s,a)\right]}_{\text{temporal-difference error}}.\]
\end{calculation}
\begin{enumerate}
\item Observe the current state $s$ and choose action $a$.
\item Receive immediate reward $r$ and new state $s'$.
\item Compute the \textbf{target} $r+\gamma\max_{a'}Q(s',a')$.
\item Move the old Q-value toward the target by fraction $\alpha$.
\end{enumerate}
\begin{keyidea}The update bootstraps: current estimates of future value improve earlier decisions.\end{keyidea}
\end{frame}

\begin{frame}{Worked Example: Update $Q((2,2),\mathrm{right})$}
\small
\conceptlead{This reproduces the actual calculation from the old lecture.}
\begin{columns}[T,onlytextwidth]
\begin{column}{.49\textwidth}
\begin{itemize}
\item Current state $(2,2)$; action \texttt{right}.
\item Next state $(2,3)$; reward $r=-10$.
\item Old estimate $Q((2,2),\mathrm{right})=5$.
\item Best Q-value in the next state is $20$.
\item Learning rate $\alpha=0.5$ and discount $\gamma=0.9$.
\end{itemize}
\end{column}
\begin{column}{.48\textwidth}
\begin{calculation}[Step 1: target]
\[y=-10+0.9(20)=8.\]
\end{calculation}
\begin{calculation}[Step 2: update]
\[Q_{\rm new}=5+0.5(8-5)=\boxed{6.5}.\]
\end{calculation}
\end{column}
\end{columns}
\lecturevisual[1.02in]{q_update_steps.png}
\end{frame}

\begin{frame}{The Learning Rate $\alpha$ and Discount Factor $\gamma$ Do Different Jobs}
\small
\begin{columns}[T,onlytextwidth]
\begin{column}{.48\textwidth}
\begin{block}{Learning rate $\alpha$}
How far should a Q-estimate move toward newly observed evidence?\\[5pt]
Small $\alpha$: gradual adjustment.\\[3pt]
Large $\alpha$: larger, potentially noisy adjustment.
\end{block}
\end{column}
\begin{column}{.48\textwidth}
\begin{block}{Discount factor $\gamma$}
How much should future rewards influence decisions today?\\[5pt]
$\gamma=0$: immediate reward only.\\[3pt]
$\gamma$ near 1: future rewards matter more.
\end{block}
\end{column}
\end{columns}
\begin{keyidea}$\alpha$ controls the \emph{update speed}; $\gamma$ controls the \emph{planning horizon}. Neither is the exploration probability.\end{keyidea}
\end{frame}

\begin{frame}{A Controlled Illustration: What $\alpha$ Changes}
\small
\conceptlead{Hold the update target fixed at $8$ and start the estimated value at $5$.}
\lecturevisual[1.86in]{alpha_effect.png}
\begin{itemize}
\item With a \emph{fixed} target, larger $\alpha$ moves more quickly toward it.
\item In a real environment, the target also changes with experience; high $\alpha$ can amplify noise.
\end{itemize}
\end{frame}

\begin{frame}{Exploration Versus Exploitation}
\small
\conceptlead{Early Q-values are uncertain; always following the current best can prevent discovery.}
\begin{columns}[T,onlytextwidth]
\begin{column}{.47\textwidth}
\begin{block}{Explore}
Try an action to collect new experience. Some actions may initially look bad only because they have not been tested.
\end{block}
\end{column}
\begin{column}{.49\textwidth}
\begin{block}{Exploit}
Choose the highest-valued available action according to the current Q-table.
\end{block}
\end{column}
\end{columns}
\begin{calculation}[$\varepsilon$-greedy strategy]
\[a=\begin{cases}
\text{random legal action},&\text{with probability }\varepsilon,\\
\arg\max_{a'}Q(s,a'),&\text{with probability }1-\varepsilon.
\end{cases}\]
\end{calculation}
\end{frame}

\begin{frame}[fragile]{An $\varepsilon$-Greedy Action Rule in Python}
\small
\conceptlead{Exploration is for training; demonstration or evaluation usually uses a greedy policy.}
{\footnotesize\begin{verbatim}
if rng.random() < epsilon:
    action = rng.choice(legal_actions)       # explore
else:
    values = Q[row, col, legal_actions]
    action = legal_actions[np.argmax(values)]  # exploit
\end{verbatim}}
\lecturevisual[1.09in]{epsilon_decay.png}
\begin{itemize}
\item Example schedule: $\varepsilon_t=\max(0.05,0.996^t)$ for episode $t$.
\item Random choices must be limited to \emph{legal} actions in our GridWorld implementation.
\end{itemize}
\end{frame}

\begin{frame}{What Is an Episode?}
\small
\conceptlead{One episode is one complete attempt between an initial state and a stopping condition.}
\begin{enumerate}
\item Reset the environment to an initial state.
\item Observe $s$, select $a$, and take a step.
\item Receive the reward and next state; update $Q(s,a)$.
\item Repeat until reaching a goal, failure condition, or step limit.
\end{enumerate}
\begin{keyidea}Learning generally takes many episodes. Episode return and task success help us tell whether behavior is improving.\end{keyidea}
\end{frame}

\section{Demo: Q-Learning on a 5 by 5 GridWorld}

\begin{frame}{Separate the Full Vacuum Task from the Teaching Demo}
\small
\conceptlead{The original lecture moves from a rich cleaning task to a simpler position-only grid.}
\begin{columns}[T,onlytextwidth]
\begin{column}{.48\textwidth}
\begin{block}{Vacuum robot (conceptual)}
State: position \emph{and} remaining dirt.\\[4pt]
Actions: move or clean.\\[4pt]
Reward: $+100$ when all cleaned; $-10$ per action.\\[4pt]
Obstacles and dirt are part of the task.
\end{block}
\end{column}
\begin{column}{.48\textwidth}
\begin{block}{GridWorld (coded demo)}
State: one of 25 positions.\\[4pt]
Actions: up, down, left, right.\\[4pt]
Reward: $+100$ on reaching goal; $-1$ otherwise.\\[4pt]
No obstacles or dirt; transitions deterministic.
\end{block}
\end{column}
\end{columns}
\begin{keyidea}Position-only Q-learning solves the smaller goal-reaching problem, not the entire vacuum-cleaning problem.\end{keyidea}
\end{frame}

\begin{frame}{GridWorld Specification for a Reproducible Run}
\small
\begin{columns}[T,onlytextwidth]
\begin{column}{.52\textwidth}
\begin{itemize}
\item A $5\times5$ grid; goal is bottom-right $(4,4)$.
\item Start each training episode at a random \emph{non-goal} position.
\item Every legal move is deterministic.
\item Entering the goal gives $+100$ and ends the episode; all other moves give $-1$.
\item Up to 100 moves per episode; a time limit is a truncation, not a goal.
\item Parameters: $\alpha=0.3$, $\gamma=0.95$; 2400 episodes; seed 4880.
\end{itemize}
\end{column}
\begin{column}{.45\textwidth}
\lecturevisual[2.58in]{gridworld_setup.png}
\end{column}
\end{columns}
\end{frame}

\begin{frame}[fragile]{Demo Setup: Q-Values and GridWorld Dynamics}
\small
\conceptlead{The Q-table has shape $5\times5\times4$ because each position has four candidate actions.}
{\footnotesize\begin{verbatim}
import numpy as np
rng = np.random.default_rng(4880)
Q = np.zeros((5, 5, 4))
ACTIONS = [(-1,0), (1,0), (0,-1), (0,1)]
GOAL = (4,4)

def step(state, action):
    row, col = state
    dr, dc = ACTIONS[action]
    nxt = (max(0, min(4, row+dr)),
           max(0, min(4, col+dc)))
    done = (nxt == GOAL)
    return nxt, (100 if done else -1), done
\end{verbatim}}
\begin{keyidea}The action-selection code filters illegal moves at boundaries.\end{keyidea}
\end{frame}

\begin{frame}[fragile]{Demo Core: Choose, Act, Update}
\small
\conceptlead{One transition is enough for one Q-learning update.}
{\footnotesize\begin{verbatim}
legal = allowed_actions(state)
if rng.random() < epsilon:
    action = int(rng.choice(legal))
else:
    action = legal[np.argmax(Q[row, col, legal])]

next_state, reward, done = step(state, action)
future = 0 if done else max(
    Q[next_state[0], next_state[1],
      allowed_actions(next_state)])

Q[row, col, action] += alpha * (
    reward + gamma * future - Q[row, col, action])
\end{verbatim}}
\begin{keyidea}For a terminal next state, the future-value contribution is zero.\end{keyidea}
\end{frame}

\begin{frame}{The Complete Training Logic at a High Level}
\small
\conceptlead{Repeat state transitions, but do not confuse training exploration with the final greedy policy.}
\lecturevisual[1.82in]{q_learning_flow.png}
\begin{itemize}
\item Each episode resets to a new starting position and uses the current $\varepsilon$.
\item After training, select the largest legal Q-value in each cell to draw the policy.
\end{itemize}
\end{frame}

\begin{frame}{Demo Results: Learning Over 2400 Episodes}
\small
\conceptlead{These supplied curves report a run using the stated seed and settings.}
\lecturevisual[2.43in]{training_curves.png}
\begin{itemize}
\item Last 100 training episodes: \textbf{100\%} goal success; mean reward \textbf{96.88}.
\item These training outcomes include $\varepsilon$-greedy exploration.
\end{itemize}
\end{frame}

\begin{frame}{The Learned Policy: Select the Largest Legal Q-Value}
\small
\begin{columns}[T,onlytextwidth]
\begin{column}{.55\textwidth}
\lecturevisual[2.33in]{learned_policy.png}
\end{column}
\begin{column}{.42\textwidth}
\begin{itemize}
\item Arrows represent $\pi(s)=\arg\max_a Q(s,a)$ after training.
\item The gold star is the terminal goal.
\item Several shortest routes can be equally good, so an alternative seed may yield different arrows.
\item This is the \textbf{simplified position-only GridWorld}.
\end{itemize}
\end{column}
\end{columns}
\end{frame}

\begin{frame}{Follow the Policy From the Top-Left Corner}
\small
\begin{columns}[T,onlytextwidth]
\begin{column}{.55\textwidth}
\lecturevisual[2.33in]{episode_trajectory.png}
\end{column}
\begin{column}{.42\textwidth}
\begin{itemize}
\item Beginning at $(0,0)$, this learned greedy rollout reaches $(4,4)$ in \textbf{8 moves}.
\item Seven moves earn $-1$ each; the final goal-reaching move earns $+100$.
\item Undiscounted episode return is \textbf{$7(-1)+100=93$}.
\item The training policy still sometimes explores; this rollout does not.
\end{itemize}
\end{column}
\end{columns}
\end{frame}

\begin{frame}{Why Not Merely Memorize the Path?}
\small
\conceptlead{Q-learning stores a decision for \emph{every} visited state, not only the one displayed route.}
\begin{itemize}
\item Random training starts create experiences from many positions.
\item The same learned table can guide the robot from alternative starting cells.
\item Bootstrapped future values propagate information about the goal back to earlier positions.
\item A grid has multiple shortest routes; the policy chooses among ties.
\item This still does \emph{not} mean the agent learned to recognize dirty tiles, plan around unseen obstacles, or clean the room.
\end{itemize}
\end{frame}

\section{Limitations and Next Steps}

\begin{frame}{Why Tabular Q-Learning Does Not Scale Well}
\small
\conceptlead{The original lecture's final slide is not a minor footnote: state-space growth is the central limitation.}
\begin{columns}[T,onlytextwidth]
\begin{column}{.51\textwidth}
\begin{itemize}
\item One value is needed per \textbf{state-action pair}.
\item A richer vacuum state must record position \emph{and} remaining dirt.
\item Learning can be slow: unvisited states receive no direct updates.
\item Continuous sensor readings and images cannot be enumerated usefully.
\item Random exploration may be costly or unsafe on real hardware.
\end{itemize}
\end{column}
\begin{column}{.46\textwidth}
\lecturevisual[2.42in]{state_space_growth.png}
\end{column}
\end{columns}
\end{frame}

\begin{frame}{From a Q-Table to a Learned Function}
\small
\conceptlead{One natural continuation is to approximate Q-values from features instead of indexing a table.}
\begin{columns}[T,onlytextwidth]
\begin{column}{.49\textwidth}
\begin{block}{Tabular Q-learning}
\[Q[s,a]\]
One independent entry for each state-action pair. Works well in small, discrete problems.
\end{block}
\end{column}
\begin{column}{.48\textwidth}
\begin{block}{Function approximation}
\[Q_{\theta}(s,a)\]
A shared model (possibly a neural network) estimates values and can generalize across similar inputs.
\end{block}
\end{column}
\end{columns}
\begin{keyidea}Function approximation helps with scale, but adds optimization and stability challenges; it is not a free fix.\end{keyidea}
\end{frame}

\begin{frame}{Common Implementation and Evaluation Mistakes}
\small
\begin{enumerate}
\item \textbf{Mixing tasks:} vacuum-with-dirt and position-only GridWorld have different state spaces and rewards.
\item \textbf{Wrong terminal update:} do not add a next-state bootstrap at a true terminal state.
\item \textbf{Unclear boundaries:} decide whether invalid moves are blocked, penalized, or excluded.
\item \textbf{Only reporting one run:} a successful seed does not prove stability across seeds.
\item \textbf{Confusing train and test behavior:} exploratory training returns can differ from greedy evaluation.
\item \textbf{Treating a learning curve as a safety claim:} toy-grid success says little about a real robot.
\end{enumerate}
\end{frame}

\begin{frame}{Connecting Reinforcement Learning Back to Intelligent Agents}
\small
\conceptlead{Reinforcement learning turns the agent loop into a method for improving behavior over time.}
\begin{itemize}
\item \textbf{Environment:} supplies a new observation and reward after each action.
\item \textbf{State:} summarizes the situation needed to choose an action.
\item \textbf{Policy:} selects a behavior in each situation.
\item \textbf{Q-values:} estimate the long-term consequences of choices.
\item \textbf{Exploration:} gathers experiences needed to improve estimates.
\item \textbf{Evaluation:} checks actual goal achievement, costs, and robustness.
\end{itemize}
\begin{keyidea}The central RL question is not ``What is the correct label?'' It is ``Which behavior leads to the best outcomes over time?''\end{keyidea}
\end{frame}

\begin{frame}{References, Provenance, and Reproducibility}
\footnotesize
\begin{itemize}
\item Source: \emph{Artificial Intelligence, Lecture 22: Learning -- Reinforcement Learning}, supplied 17-page slide deck, especially pages 2--17 (agent, vacuum grid, states, rewards, policies, Q-tables, update, worked example, $\varepsilon$-greedy, episodes, code, and limitations).
\item Reproduced old-lecture worked result: $5+0.5[-10+0.9(20)-5]=6.5$.
\item Demo extension: the supplied position-only $5\times5$ Q-learning experiment uses explicit legal actions, terminal handling, a fixed seed, and separate exploration/evaluation.
\item Diagrams and reported GridWorld outcomes are included in the supplied image archive. The original figure-generation script was not included, so the reported training statistics have not been independently reproduced.
\item Further reading: Sutton and Barto, \emph{Reinforcement Learning: An Introduction} (2nd ed., 2018), Chapters 3 and 6; Watkins and Dayan (1992), \emph{Q-learning}.
\item Measured results may vary with changes to code, RNG, or numerical dependencies; the supplied run uses NumPy's \texttt{default\_rng(4880)}.
\end{itemize}
\end{frame}
\end{document}
